\documentclass[11pt,a4paper]{article}
\usepackage[utf8]{inputenc}
\usepackage[T1]{fontenc}
\usepackage[spanish,es-minimal,es-noshorthands]{babel}
\usepackage{lmodern}
\usepackage[sfdefault,lining]{FiraSans}
\usepackage[a4paper,top=2.3cm,bottom=2.2cm,left=1.8cm,right=1.8cm,headheight=1.3cm,headsep=0.4cm,footskip=1cm]{geometry}
\usepackage{amsmath,amssymb}
\usepackage{icomma}
\usepackage{graphicx}
\usepackage[table]{xcolor}
\usepackage{tikz}
\usetikzlibrary{calc}
\usepackage{fontawesome5}
\usepackage{ccicons}
\usepackage{fancyhdr}
\usepackage{lastpage}
\usepackage{tasks}
\usepackage{needspace}
\usepackage[most]{tcolorbox}
\usepackage{array}
\usepackage[hidelinks]{hyperref}
\AtBeginDocument{\hypersetup{pdftitle={Soluciones de la ficha de la unidad 2: Álgebra. Ecuaciones e inecuaciones. 1º Bach.},pdfsubject={Matemáticas I 1º Bach.},pdfkeywords={polinomios, Ruffini, teorema del resto, factorización, fracciones algebraicas, ecuaciones, inecuaciones}}}

\definecolor{navy}{HTML}{1B2748}
\definecolor{accent}{HTML}{E8603A}
\definecolor{teal}{HTML}{1F9E90}
\definecolor{soft}{HTML}{EEF1F8}
\definecolor{muted}{HTML}{5D6685}

\setlength{\parindent}{0pt}
\setlength{\parskip}{3pt}
\settasks{label=\alph*), label-format=\bfseries\color{navy}, label-width=1.3em, item-indent=1.8em, column-sep=0.6em, after-item-skip=5pt, before-skip=3pt, after-skip=4pt}

% ---------- capçalera i peu ----------
\newcommand{\llicencia}{\href{https://creativecommons.org/licenses/by-nc-sa/4.0/deed.es}{\ccbyncsa}\ \ Licencia Creative Commons Atribución-NoComercial-CompartirIgual 4.0 Internacional (CC BY-NC-SA 4.0)}
\pagestyle{fancy}
\fancyhf{}
\fancyhead[L]{\raisebox{-0.35\height}{\includegraphics[height=0.95cm]{logo-gva.png}}}
\fancyhead[C]{\small\color{muted}Soluciones de la ficha · Unidad 2: Álgebra. Ecuaciones e inecuaciones · 1\textsuperscript{o} Bach.}
\fancyhead[R]{\raisebox{-0.35\height}{\includegraphics[height=0.95cm]{logo-vermellar.png}}}
\fancyfoot[L]{\scriptsize\color{muted}\llicencia}
\fancyfoot[R]{\small\color{muted}\thepage}
\renewcommand{\headrulewidth}{0.4pt}
\renewcommand{\headrule}{\color{muted!40}\hrule height \headrulewidth}
\fancypagestyle{primera}{\fancyhf{}\renewcommand{\headrulewidth}{0pt}
  \fancyfoot[L]{\scriptsize\color{muted}\llicencia}
  \fancyfoot[R]{\small\color{muted}\thepage}}

% ---------- nivells i exercicis ----------
\newcommand{\estrelles}[1]{{\color{accent}\ifcase#1\or\faStar\faStar[regular]\faStar[regular]\or\faStar\faStar\faStar[regular]\or\faStar\faStar\faStar\fi}}
\newcounter{exer}
\newcommand{\ex}[1]{\par\needspace{4\baselineskip}\medskip\refstepcounter{exer}\noindent{\color{navy}\bfseries\theexer.}\ \ }

\newcommand{\seccio}[2]{\par\needspace{16\baselineskip}\bigskip
  \begin{tcolorbox}[colback=navy,colframe=navy,coltext=white,arc=2mm,boxrule=0pt,left=3mm,right=3mm,top=1.2mm,bottom=1.2mm]
  {\large\bfseries #1\quad #2}
  \end{tcolorbox}\vspace{-2pt}}

\newcommand{\exf}[1]{\par\medskip\refstepcounter{exer}\noindent{\color{navy}\bfseries\theexer.}\ \ }
\providecommand{\lt}{<}\providecommand{\gt}{>}
\newcolumntype{C}[1]{>{\centering\arraybackslash}p{#1}}
% quadre de recordatori (només a la fitxa d'exercicis)
\newcommand{\recorda}[1]{\par\smallskip\begin{tcolorbox}[colback=accent!7,colframe=accent!55,arc=2mm,boxrule=0.5pt,left=2.5mm,right=2.5mm,top=1mm,bottom=1mm]\small{\color{accent}\faLightbulb[regular]}\ \textbf{Recuerda:}\ #1\end{tcolorbox}}
\newcommand{\res}[1]{\boxed{#1}}\newcommand{\rest}[1]{\fbox{#1}}

\def\fmneed{11}% que el títol d'una secció no quede sol al final del full
% =====================================================================
% Capçalera FentMates de les fitxes: un color per a cada curs.
% Els colors són els de la pàgina inicial (LEVELS a assets/js/materials.js);
% el logo ve de _eines/logo/pdf/ (python3 _eines/logo/genera.py).
%
% Ús, al final del preàmbul de la unitat:
%   % =====================================================================
% Capçalera FentMates de les fitxes: un color per a cada curs.
% Els colors són els de la pàgina inicial (LEVELS a assets/js/materials.js);
% el logo ve de _eines/logo/pdf/ (python3 _eines/logo/genera.py).
%
% Ús, al final del preàmbul de la unitat:
%   % =====================================================================
% Capçalera FentMates de les fitxes: un color per a cada curs.
% Els colors són els de la pàgina inicial (LEVELS a assets/js/materials.js);
% el logo ve de _eines/logo/pdf/ (python3 _eines/logo/genera.py).
%
% Ús, al final del preàmbul de la unitat:
%   \input{../comu/fentmates.tex}
%   \fmcurs{1eso}
% i, al començament del document:
%   \capcalera{Matemàtiques}{1r ESO}{Fitxa d'exercicis · Unitat 2}{Divisibilitat}
% Opcional abans de \input: \def\fmneed{14} (espai mínim abans de cada secció).
% =====================================================================
\makeatletter
\def\fm@logodir{../../logo/pdf}
\providecommand{\fmneed}{10}
\def\fm@cursnom{}

% clau del curs, color i color del text damunt del color
\newcommand{\fm@curs}[3]{\@namedef{fm@color@#1}{#2}\@namedef{fm@text@#1}{#3}}
\fm@curs{1eso}{E6B84A}{navy}
\fm@curs{2eso}{3B82F6}{white}
\fm@curs{3eso}{FF7A50}{white}
\fm@curs{4eso-a}{A78BFA}{white}
\fm@curs{4eso-b}{6D28D9}{white}
\fm@curs{1batx-mi}{9B1C1C}{white}
\fm@curs{1batx-ccss}{EF6B6B}{white}
\fm@curs{2batx-mii}{1E6B47}{white}
\fm@curs{2batx-ccss}{2BB3A3}{white}

\newcommand{\fmcurs}[1]{%
  \@ifundefined{fm@color@#1}{\PackageError{fentmates}{Curs desconegut: #1}{}}{}%
  \edef\fm@tmp{\noexpand\definecolor{curs}{HTML}{\@nameuse{fm@color@#1}}%
               \noexpand\colorlet{curstext}{\@nameuse{fm@text@#1}}}\fm@tmp
  \colorlet{cursfosc}{curs!72!black}%
  \def\fm@logo{\fm@logodir/vertical-#1.pdf}%
  \def\fm@marca{\fm@logodir/marca-#1.pdf}}

% ---------- capçalera de la primera pàgina ----------
% #1 matèria (petit), #2 curs, #3 tipus de document i unitat, #4 títol
\newcommand{\capcalera}[4]{%
  \gdef\fm@cursnom{#2}%
  \par\noindent
  \begin{tcolorbox}[colback=curs,colframe=curs,boxrule=0pt,arc=2.2mm,width=3.5cm,height=2.15cm,
      valign=center,left=3mm,right=2mm,top=1mm,bottom=1mm,nobeforeafter,box align=center]
    \color{curstext}{\footnotesize #1}\par\vspace{1pt}{\fontsize{19}{21}\selectfont\bfseries #2\par}%
  \end{tcolorbox}\hspace{5mm}%
  \begin{minipage}[c]{\dimexpr\linewidth-3.5cm-5mm-2.9cm\relax}
    {\color{cursfosc}\small\bfseries\MakeUppercase{#3}}\par\vspace{2pt}
    {\color{navy}\fontsize{20}{23}\selectfont\bfseries #4\par}
  \end{minipage}\hfill
  \begin{minipage}[c]{2.6cm}\raggedleft\includegraphics[width=2.6cm]{\fm@logo}\end{minipage}\par
  \vspace{2pt}}

% ---------- barres de secció del color del curs ----------
\renewcommand{\seccio}[2]{\par\needspace{\fmneed\baselineskip}\bigskip
  \begin{tcolorbox}[colback=curs,colframe=curs,coltext=curstext,arc=2mm,boxrule=0pt,left=3mm,right=3mm,top=1.2mm,bottom=1.2mm]
  {\large\bfseries #1\quad #2}
  \end{tcolorbox}\vspace{-2pt}}

% ---------- estrelles de dificultat del color del curs (a Batxillerat no se'n posen) ----------
\@ifundefined{estrelles}{}{%
  \renewcommand{\estrelles}[1]{{\color{curs}\ifcase#1\or\faStar\faStar[regular]\faStar[regular]\or\faStar\faStar\faStar[regular]\or\faStar\faStar\faStar\fi}}}

% ---------- capçalera i peu de les pàgines ----------
% primera pàgina: logos institucionals al peu
\fancypagestyle{primera}{\fancyhf{}\renewcommand{\headrulewidth}{0pt}
  \fancyfoot[L]{\raisebox{-0.32\height}{\includegraphics[height=0.78cm]{logo-gva.png}}\hspace{5mm}%
    \raisebox{-0.32\height}{\includegraphics[height=0.78cm]{logo-vermellar.png}}}
  \fancyfoot[R]{\scriptsize\color{muted}\href{https://creativecommons.org/licenses/by-nc-sa/4.0/deed.ca}{\ccbyncsa}\ \ CC BY-NC-SA 4.0 · fentmates.org\qquad\small\thepage}}
% la resta: el símbol a l'esquerra i el curs a la dreta, amb el títol al centre com abans
\fancyhead[L]{\raisebox{-0.3\height}{\includegraphics[height=0.62cm]{\fm@marca}}}
\fancyhead[R]{\small\bfseries\color{cursfosc}\fm@cursnom}
\renewcommand{\headrule}{\color{curs!45}\hrule height \headrulewidth}
\makeatother

%   \fmcurs{1eso}
% i, al començament del document:
%   \capcalera{Matemàtiques}{1r ESO}{Fitxa d'exercicis · Unitat 2}{Divisibilitat}
% Opcional abans de \input: \def\fmneed{14} (espai mínim abans de cada secció).
% =====================================================================
\makeatletter
\def\fm@logodir{../../logo/pdf}
\providecommand{\fmneed}{10}
\def\fm@cursnom{}

% clau del curs, color i color del text damunt del color
\newcommand{\fm@curs}[3]{\@namedef{fm@color@#1}{#2}\@namedef{fm@text@#1}{#3}}
\fm@curs{1eso}{E6B84A}{navy}
\fm@curs{2eso}{3B82F6}{white}
\fm@curs{3eso}{FF7A50}{white}
\fm@curs{4eso-a}{A78BFA}{white}
\fm@curs{4eso-b}{6D28D9}{white}
\fm@curs{1batx-mi}{9B1C1C}{white}
\fm@curs{1batx-ccss}{EF6B6B}{white}
\fm@curs{2batx-mii}{1E6B47}{white}
\fm@curs{2batx-ccss}{2BB3A3}{white}

\newcommand{\fmcurs}[1]{%
  \@ifundefined{fm@color@#1}{\PackageError{fentmates}{Curs desconegut: #1}{}}{}%
  \edef\fm@tmp{\noexpand\definecolor{curs}{HTML}{\@nameuse{fm@color@#1}}%
               \noexpand\colorlet{curstext}{\@nameuse{fm@text@#1}}}\fm@tmp
  \colorlet{cursfosc}{curs!72!black}%
  \def\fm@logo{\fm@logodir/vertical-#1.pdf}%
  \def\fm@marca{\fm@logodir/marca-#1.pdf}}

% ---------- capçalera de la primera pàgina ----------
% #1 matèria (petit), #2 curs, #3 tipus de document i unitat, #4 títol
\newcommand{\capcalera}[4]{%
  \gdef\fm@cursnom{#2}%
  \par\noindent
  \begin{tcolorbox}[colback=curs,colframe=curs,boxrule=0pt,arc=2.2mm,width=3.5cm,height=2.15cm,
      valign=center,left=3mm,right=2mm,top=1mm,bottom=1mm,nobeforeafter,box align=center]
    \color{curstext}{\footnotesize #1}\par\vspace{1pt}{\fontsize{19}{21}\selectfont\bfseries #2\par}%
  \end{tcolorbox}\hspace{5mm}%
  \begin{minipage}[c]{\dimexpr\linewidth-3.5cm-5mm-2.9cm\relax}
    {\color{cursfosc}\small\bfseries\MakeUppercase{#3}}\par\vspace{2pt}
    {\color{navy}\fontsize{20}{23}\selectfont\bfseries #4\par}
  \end{minipage}\hfill
  \begin{minipage}[c]{2.6cm}\raggedleft\includegraphics[width=2.6cm]{\fm@logo}\end{minipage}\par
  \vspace{2pt}}

% ---------- barres de secció del color del curs ----------
\renewcommand{\seccio}[2]{\par\needspace{\fmneed\baselineskip}\bigskip
  \begin{tcolorbox}[colback=curs,colframe=curs,coltext=curstext,arc=2mm,boxrule=0pt,left=3mm,right=3mm,top=1.2mm,bottom=1.2mm]
  {\large\bfseries #1\quad #2}
  \end{tcolorbox}\vspace{-2pt}}

% ---------- estrelles de dificultat del color del curs (a Batxillerat no se'n posen) ----------
\@ifundefined{estrelles}{}{%
  \renewcommand{\estrelles}[1]{{\color{curs}\ifcase#1\or\faStar\faStar[regular]\faStar[regular]\or\faStar\faStar\faStar[regular]\or\faStar\faStar\faStar\fi}}}

% ---------- capçalera i peu de les pàgines ----------
% primera pàgina: logos institucionals al peu
\fancypagestyle{primera}{\fancyhf{}\renewcommand{\headrulewidth}{0pt}
  \fancyfoot[L]{\raisebox{-0.32\height}{\includegraphics[height=0.78cm]{logo-gva.png}}\hspace{5mm}%
    \raisebox{-0.32\height}{\includegraphics[height=0.78cm]{logo-vermellar.png}}}
  \fancyfoot[R]{\scriptsize\color{muted}\href{https://creativecommons.org/licenses/by-nc-sa/4.0/deed.ca}{\ccbyncsa}\ \ CC BY-NC-SA 4.0 · fentmates.org\qquad\small\thepage}}
% la resta: el símbol a l'esquerra i el curs a la dreta, amb el títol al centre com abans
\fancyhead[L]{\raisebox{-0.3\height}{\includegraphics[height=0.62cm]{\fm@marca}}}
\fancyhead[R]{\small\bfseries\color{cursfosc}\fm@cursnom}
\renewcommand{\headrule}{\color{curs!45}\hrule height \headrulewidth}
\makeatother

%   \fmcurs{1eso}
% i, al començament del document:
%   \capcalera{Matemàtiques}{1r ESO}{Fitxa d'exercicis · Unitat 2}{Divisibilitat}
% Opcional abans de \input: \def\fmneed{14} (espai mínim abans de cada secció).
% =====================================================================
\makeatletter
\def\fm@logodir{../../logo/pdf}
\providecommand{\fmneed}{10}
\def\fm@cursnom{}

% clau del curs, color i color del text damunt del color
\newcommand{\fm@curs}[3]{\@namedef{fm@color@#1}{#2}\@namedef{fm@text@#1}{#3}}
\fm@curs{1eso}{E6B84A}{navy}
\fm@curs{2eso}{3B82F6}{white}
\fm@curs{3eso}{FF7A50}{white}
\fm@curs{4eso-a}{A78BFA}{white}
\fm@curs{4eso-b}{6D28D9}{white}
\fm@curs{1batx-mi}{9B1C1C}{white}
\fm@curs{1batx-ccss}{EF6B6B}{white}
\fm@curs{2batx-mii}{1E6B47}{white}
\fm@curs{2batx-ccss}{2BB3A3}{white}

\newcommand{\fmcurs}[1]{%
  \@ifundefined{fm@color@#1}{\PackageError{fentmates}{Curs desconegut: #1}{}}{}%
  \edef\fm@tmp{\noexpand\definecolor{curs}{HTML}{\@nameuse{fm@color@#1}}%
               \noexpand\colorlet{curstext}{\@nameuse{fm@text@#1}}}\fm@tmp
  \colorlet{cursfosc}{curs!72!black}%
  \def\fm@logo{\fm@logodir/vertical-#1.pdf}%
  \def\fm@marca{\fm@logodir/marca-#1.pdf}}

% ---------- capçalera de la primera pàgina ----------
% #1 matèria (petit), #2 curs, #3 tipus de document i unitat, #4 títol
\newcommand{\capcalera}[4]{%
  \gdef\fm@cursnom{#2}%
  \par\noindent
  \begin{tcolorbox}[colback=curs,colframe=curs,boxrule=0pt,arc=2.2mm,width=3.5cm,height=2.15cm,
      valign=center,left=3mm,right=2mm,top=1mm,bottom=1mm,nobeforeafter,box align=center]
    \color{curstext}{\footnotesize #1}\par\vspace{1pt}{\fontsize{19}{21}\selectfont\bfseries #2\par}%
  \end{tcolorbox}\hspace{5mm}%
  \begin{minipage}[c]{\dimexpr\linewidth-3.5cm-5mm-2.9cm\relax}
    {\color{cursfosc}\small\bfseries\MakeUppercase{#3}}\par\vspace{2pt}
    {\color{navy}\fontsize{20}{23}\selectfont\bfseries #4\par}
  \end{minipage}\hfill
  \begin{minipage}[c]{2.6cm}\raggedleft\includegraphics[width=2.6cm]{\fm@logo}\end{minipage}\par
  \vspace{2pt}}

% ---------- barres de secció del color del curs ----------
\renewcommand{\seccio}[2]{\par\needspace{\fmneed\baselineskip}\bigskip
  \begin{tcolorbox}[colback=curs,colframe=curs,coltext=curstext,arc=2mm,boxrule=0pt,left=3mm,right=3mm,top=1.2mm,bottom=1.2mm]
  {\large\bfseries #1\quad #2}
  \end{tcolorbox}\vspace{-2pt}}

% ---------- estrelles de dificultat del color del curs (a Batxillerat no se'n posen) ----------
\@ifundefined{estrelles}{}{%
  \renewcommand{\estrelles}[1]{{\color{curs}\ifcase#1\or\faStar\faStar[regular]\faStar[regular]\or\faStar\faStar\faStar[regular]\or\faStar\faStar\faStar\fi}}}

% ---------- capçalera i peu de les pàgines ----------
% primera pàgina: logos institucionals al peu
\fancypagestyle{primera}{\fancyhf{}\renewcommand{\headrulewidth}{0pt}
  \fancyfoot[L]{\raisebox{-0.32\height}{\includegraphics[height=0.78cm]{logo-gva.png}}\hspace{5mm}%
    \raisebox{-0.32\height}{\includegraphics[height=0.78cm]{logo-vermellar.png}}}
  \fancyfoot[R]{\scriptsize\color{muted}\href{https://creativecommons.org/licenses/by-nc-sa/4.0/deed.ca}{\ccbyncsa}\ \ CC BY-NC-SA 4.0 · fentmates.org\qquad\small\thepage}}
% la resta: el símbol a l'esquerra i el curs a la dreta, amb el títol al centre com abans
\fancyhead[L]{\raisebox{-0.3\height}{\includegraphics[height=0.62cm]{\fm@marca}}}
\fancyhead[R]{\small\bfseries\color{cursfosc}\fm@cursnom}
\renewcommand{\headrule}{\color{curs!45}\hrule height \headrulewidth}
\makeatother

\fmcurs{1batx-mi}

\begin{document}
\thispagestyle{primera}

\capcalera{Matemáticas I}{1\textsuperscript{o} Bach.}{Soluciones de la ficha · Unidad 2}{\raggedright Álgebra. Ecuaciones e inecuaciones}
\vspace{8pt}

\begin{tcolorbox}[colback=soft,colframe=soft,arc=2mm,left=3mm,right=3mm,top=1.5mm,bottom=1.5mm]
\small Intenta resolver cada ejercicio antes de mirar su solución. Los resultados se dan de manera exacta (fracciones y raíces) y, cuando hace falta, con el valor aproximado. Puedes comprobar las divisiones, la factorización y las inecuaciones con las herramientas interactivas de los apuntes de la unidad.
\end{tcolorbox}

% =====================================================
\seccio{1}{Polinomios. Operaciones}

\exf{1} \begin{tasks}(1)
  \task $\res{- 16 x^{2} - 17 x + 10}$
  \task $\res{15 x^{2} + 3 x - 3}$
  \task $\res{6 x^{4} - 10 x^{3} - 35 x^{2} + 23 x}$
  \task $\res{3 x^{5} - 2 x^{4} - 20 x^{2} - 16 x + 32}$
  \task $\res{x^{3} - \frac{9 x^{2}}{10} + \frac{4 x}{15}}$
\end{tasks}

\ex{2} $(x+y+z)^2=\big((x+y)+z\big)^2=(x+y)^2+2(x+y)z+z^2=x^2+2xy+y^2+2xz+2yz+z^2$. \rest{Queda demostrada.}

\ex{1} $\res{16 x^{4} - 24 x^{3} + 25 x^{2} - 12 x + 4}$ \ (con la identidad del ejercicio anterior).

\ex{2} \begin{tasks}(1)
  \task $6x^4+(3b-10)x^3+(2a-5b)x^2+abx$: \ $3b-10=-19$, $ab=-6\Rightarrow \res{a=2,\ b=-3}$
  \task $2x^2+9x+18-3a$: \ $\res{b=9},\ \ 18-3a=3\Rightarrow \res{a=5}$
\end{tasks}

\ex{1} \begin{tasks}(1)
  \task $c(x)=\res{x + \frac{2}{3}},\ \ r(x)=\res{- x - \frac{19}{3}}$
  \task $c(x)=\res{3 x^{2} - 3 x - 5},\ \ r(x)=\res{10 x + 14}$
  \task $c(x)=\res{\frac{1}{2}},\ \ r(x)=\res{- 3 x^{2} + \frac{9}{2}}$
\end{tasks}

\ex{1} \begin{tasks}(1)
  \task ${\def\arraystretch{1.5}\begin{array}{r|rrr|r} & 3 & 2 & 1 & -5\\-2 &  & -6 & 8 & -18\\\hline & 3 & -4 & 9 & -23\end{array}}$\\[2pt]$c(x)=\res{3 x^{2} - 4 x + 9},\ \ r=\res{-23}$
  \task ${\def\arraystretch{1.5}\begin{array}{r|rrrr|r} & 3 & 0 & -2 & -1 & 4\\-2 &  & -6 & 12 & -20 & 42\\\hline & 3 & -6 & 10 & -21 & 46\end{array}}$\\[2pt]$c(x)=\res{3 x^{3} - 6 x^{2} + 10 x - 21},\ \ r=\res{46}$
  \task ${\def\arraystretch{1.5}\begin{array}{r|rrr|r} & 1 & -3 & -1 & 6\\- \frac{3}{2} &  & - \frac{3}{2} & \frac{27}{4} & - \frac{69}{8}\\\hline & 1 & - \frac{9}{2} & \frac{23}{4} & - \frac{21}{8}\end{array}}$\\[2pt]$c(x)=\res{\frac{x^{2}}{2} - \frac{9 x}{4} + \frac{23}{8}},\ \ r=\res{- \frac{21}{8}}\ \ \text{(cociente de Ruffini dividido entre }2)$
  \task ${\def\arraystretch{1.5}\begin{array}{r|rrrr|r} & 4 & 0 & -3 & 0 & 5\\5 &  & 20 & 100 & 485 & 2425\\\hline & 4 & 20 & 97 & 485 & 2430\end{array}}$\\[2pt]$c(x)=\res{- 4 x^{3} - 20 x^{2} - 97 x - 485},\ \ r=\res{2430}\ \ \text{(cociente de Ruffini dividido entre }-1)$
  \task ${\def\arraystretch{1.5}\begin{array}{r|rrrrr|r} & 2 & 0 & -12 & -6 & 0 & -117\\3 &  & 6 & 18 & 18 & 36 & 108\\\hline & 2 & 6 & 6 & 12 & 36 & -9\end{array}}$\\[2pt]$c(x)=\res{2 x^{4} + 6 x^{3} + 6 x^{2} + 12 x + 36},\ \ r=\res{-9}$
\end{tasks}

\ex{2} \begin{tasks}(1)
  \task $P(-3)=a - 3=0\Rightarrow \res{a=3}$
  \task $P(- \frac{3}{2})=\frac{3 a}{2} - \frac{3}{2}=0\Rightarrow \res{a=1}$
  \task $P(\frac{4}{3})=\frac{16 a}{9} + \frac{160}{9}=0\Rightarrow \res{a=-10}$
  \task $P(5)=b + 2930=0\Rightarrow \res{b=-2930}$
\end{tasks}(Con el teorema del resto, $P(r)$ debe ser 0, donde $r$ es la raíz del divisor; es lo mismo que obtener resto 0 en Ruffini.)

\ex{1} $M(2)=16-20+8-4=0$: \rest{sí es divisible}.\par Por Ruffini: ${\def\arraystretch{1.5}\begin{array}{r|rrr|r} & 2 & -5 & 4 & -4\\2 &  & 4 & -2 & 4\\\hline & 2 & -1 & 2 & 0\end{array}}\Rightarrow \res{N(x)=2x^2-x+2}$


% =====================================================
\seccio{2}{Factorización de un polinomio. Teorema del resto}

\exf{1} \begin{tasks}(5)
  \task $P(1)=0$: \rest{sí}
  \task $P(2)=36$: \rest{no}
  \task $P(-1)=-18$: \rest{no}
  \task $P(-4)=0$: \rest{sí}
  \task $P(-3)=-44$: \rest{no}
\end{tasks}

\ex{1} $P(-1)=1+2-3=\res{0}$, \ $P(3)=9-6-3=\res{0}$. \rest{No}: un polinomio de grado 2 tiene como máximo dos raíces, y ya las hemos encontrado: $P(x)=(x+1)(x-3)$.

\ex{1} \begin{tasks}(1)
  \task $\res{x=-5,\ x=2,\ x=3}$
  \task $\res{x=- \frac{1}{2},\ x=0,\ x=2\ (\text{triple})}$
  \task $\res{x=-1\ (\text{triple}),\ x=- \frac{2}{3},\ x=\frac{1}{2}}$
  \task $\res{x=-2,\ x=1,\ x=4}$ \ $\big((x - 4) (x - 1) (x + 2)\big)$
  \task $\res{x=-5,\ x=-2,\ x=-1}$ \ $\big((x + 1) (x + 2) (x + 5)\big)$
  \task $\res{x=0,\ x=\frac{1}{2},\ x=1,\ x=2\ (\text{doble})}$ \ $\big(x (x - 2)^{2} (x - 1) (2 x - 1)\big)$
\end{tasks}

\ex{2} \begin{tasks}(1)
  \task $\res{(x - 3) (x - 1) (x + 2)}$ \ \small(raíces: $x=-2,\ x=1,\ x=3$)
  \task $\res{(x - 1)^{2} (x + 3)}$ \ \small(raíces: $x=-3,\ x=1\ (\text{doble})$)
  \task $\res{(x - 1) (x + 1)^{2} (2 x - 3)}$ \ \small(raíces: $x=-1\ (\text{doble}),\ x=1,\ x=\frac{3}{2}$)
  \task $\res{(x - 2) (2 x - 1) (2 x + 3)}$ \ \small(raíces: $x=- \frac{3}{2},\ x=\frac{1}{2},\ x=2$)
  \task $\res{x^{2} (x - 2) (x + 2)}$ \ \small(raíces: $x=-2,\ x=0\ (\text{doble}),\ x=2$)
  \task $\res{x (x + 1)^{2} (2 x - 1) (3 x + 1)}$ \ \small(raíces: $x=-1\ (\text{doble}),\ x=- \frac{1}{3},\ x=0,\ x=\frac{1}{2}$)
  \task $\res{(x - 2) (x - 1) (x + 1) (x + 2) (x^{2} + 1)}$ \ \small(raíces: $x=-2,\ x=-1,\ x=1,\ x=2$)
  \task $t=x^2$: $\frac{t^2}{4}+\frac34t+\frac9{16}=0\Rightarrow\Delta=\frac9{16}-\frac9{16}=0$, $t=-\frac32\lt0$: \rest{no tiene raíces reales}. $P(x)=\res{\frac{1}{16}(2x^2+3)^2}$
\end{tasks}

\ex{1} \begin{tasks}(1)
  \task $P(2)=k + 8=0\Rightarrow \res{k=-8}$
  \task $P(-2)=- 2 k - 12=0\Rightarrow \res{k=-6}$
  \task $P(-4)=- 63 k + 252=0\Rightarrow \res{k=4}$
  \task $P(3)=- 6 k + 36=0\Rightarrow \res{k=6}$
\end{tasks}

\ex{2} $P(-3)=-27-9(3k+2)-3=-48-27k$.\begin{tasks}(2)
  \task $-48-27k=0\Rightarrow \res{k=-\frac{16}{9}}$
  \task $-48-27k=15\Rightarrow \res{k=-\frac73}$
\end{tasks}

\ex{2} $P(x)=ax^2+bx+c$: \ $c=-3$, \ $a+b-3=-6$, \ $9a+3b-3=0\Rightarrow a=2,\ b=-5$: \ $\res{P(x)=2x^2-5x-3}$

\ex{1} $P(2)=8m-24-8+8=8m-24=0\Rightarrow \res{m=3}$

\ex{2} $P(2)=0$: $4a+2b=-2$; \ $P(-3)=0$: $9a-3b=33$. \ $\res{a=2,\ b=-5}$ \ (y entonces $P(x)=(x-2)(x+3)(x+1)$).

\ex{1} Teorema del resto: $P(-2)=32-72+8+12+3m=-20+3m=12\Rightarrow \res{m=\frac{32}{3}}$


% =====================================================
\seccio{3}{Fracciones algebraicas}

\exf{2} \begin{tasks}(1)
  \task $\dfrac{2 (x - 2) (x + 2)}{x (x - 1) (x + 2)}=\res{\frac{2 (x - 2)}{x (x - 1)}}$
  \task $\dfrac{x (x - 1) (x + 1)}{x (x + 1) (x + 2)}=\res{\frac{x - 1}{x + 2}}$
  \task $\dfrac{x^{2} (x - 3) (x + 1)}{x^{2} (x - 3) (x + 3)}=\res{\frac{x + 1}{x + 3}}$
  \task $\dfrac{(x - 1)^{3}}{x (x - 1)^{2}}=\res{\frac{x - 1}{x}}$
  \task $\dfrac{(x - 1)^{2} (x + 1) (2 x - 3)}{(x - 1) (x + 2) (2 x - 3)}=\res{\frac{(x - 1) (x + 1)}{x + 2}}$
  \task $\dfrac{(x - 1)^{2} (x + 3) (2 x - 1)}{(x - 1) (x + 3) (2 x - 1)}=\res{x - 1}$
  \task $\dfrac{(x - 3) (x - 2) (x + 3)}{(x - 3) (x - 2)^{2}}=\res{\frac{x + 3}{x - 2}}$
  \task $\dfrac{(x - 2) (x - 1)^{3}}{(x - 2)^{2} (x - 1)^{2}}=\res{\frac{x - 1}{x - 2}}$
\end{tasks}

\ex{2} \begin{tasks}(2)
  \task $\res{\frac{x + 2}{(x - 1) (x + 1)}}$
  \task $\res{\frac{3 x + 1}{x (x + 1)}}$
  \task $\res{\frac{2 x - 5}{4 x (x - 1)}}$
  \task $\res{\frac{2 (2 x + 1)}{x (x - 1) (x + 1)}}$
  \task $\res{\frac{5 x^{2} - 3 x + 1}{(x - 2) (x + 2)}}$
  \task $\res{\frac{2 (x^{2} + x - 1)}{(x - 1)^{2} (x + 1)}}$
\end{tasks}

\ex{2} \begin{tasks}(2)
  \task $\res{\frac{1}{4}}$
  \task $\res{\frac{x + 3}{(x - 2) (x + 1)}}$
  \task $\res{\frac{(3 x + 2) (4 x + 1)}{(x - 2) (x + 2)}}$
  \task $\res{- \frac{2}{3 x^{2}}}$
  \task $\res{- \frac{2 x^{2} - 3 x - 3}{x - 1}}$
  \task $\res{\frac{x^{2} - 6 x - 1}{2 (x + 1)^{2}}}$
  \task $\res{x}$
\end{tasks}

\ex{3} \begin{tasks}(2)
  \task $\res{\frac{3 x^{4} + 3 x^{3} + 3 x^{2} + 4 x + 2}{x (x + 2) (x^{2} + 1)}}$
  \task $\res{\frac{3 x (x + 1)}{x^{2} + 1}}$
  \task $\res{- \frac{a + b}{b^{2}}}$
  \task $\res{- \frac{a (b^{2} - b - 1)}{b^{2}}}$
  \task $\res{\frac{2 (5 x + 6)}{(x - 2) (x + 2)}}$
  \task $\res{\frac{(x - 2) (x + 5)}{x + 1}}$
  \task $\res{x y (x - y)}$
  \task $\res{\frac{1}{a + x}}$
  \task $\res{\frac{3 x + 2}{2 x + 1}}$
  \task $\res{\frac{3 x^{2} + 3 x + 2}{x^{2} + 1}}$
\end{tasks}


% =====================================================
\seccio{4}{Ecuaciones polinómicas}

\exf{2} \begin{tasks}(1)
  \task $\res{x=\frac{9}{32}}$
  \task $\res{x=-5,\ x=0}$
  \task $\res{x=- \frac{\sqrt{30}}{6},\ x=\frac{\sqrt{30}}{6}}$
  \task $\res{x=-2,\ x=\frac{11}{4}}$
  \task $(x - 1) (x + 2) (2 x + 3) (3 x - 1)=0$: \ $\res{x=-2,\ x=- \frac{3}{2},\ x=\frac{1}{3},\ x=1}$
  \task $(x - 2) (2 x + 3) (x^{2} + 3)=0$: \ $\res{x=- \frac{3}{2},\ x=2}$
  \task $x (x + 1) (2 x - 1) (3 x + 1)=0$: \ $\res{x=-1,\ x=- \frac{1}{3},\ x=0,\ x=\frac{1}{2}}$
  \task $(x - 2) (x - 1) (x + 3)=0$: \ $\res{x=-3,\ x=1,\ x=2}$
  \task $(x - 3)^{2} (x + 3)=0$: \ $\res{x=-3,\ x=3}$
  \task $(x - 3) (x + 2) (x^{2} + 1)=0$: \ $\res{x=-2,\ x=3}$
  \task $(x - 1) (x + 1) (x^{2} + 1)=0$: \ $\res{x=-1,\ x=1}$
  \task $x (x - 3) (x + 1) (x + 3)=0$: \ $\res{x=-3,\ x=-1,\ x=0,\ x=3}$
\end{tasks}

\ex{2} \begin{tasks}(1)
  \task $t=x^2$: $t^2-17t+16=0\Rightarrow t=16,\ t=1$: \ $\res{x=\pm4,\ x=\pm1}$
  \task $t=x^3$: $t^2-9t+8=0\Rightarrow t=8,\ t=1$: \ $\res{x=2,\ x=1}$
  \task $t^2-8t-9=0\Rightarrow t=9$ ($t=-1$ no da soluciones): \ $\res{x=\pm3}$
  \task $t=x^3$: $t^2+2t-3=0\Rightarrow t=1,\ t=-3$: \ $\res{x=1,\ x=-\sqrt[3]{3}}$
  \task Desarrollando: $2x^4+12x^2-14=0\Rightarrow t^2+6t-7=0\Rightarrow t=1$ ($t=-7$ no da soluciones): \ $\res{x=\pm1}$
  \task Multiplicamos por $x^2$: $x^4-10x^2+9=0\Rightarrow t=9,\ t=1$: \ $\res{x=\pm3,\ x=\pm1}$
\end{tasks}

\ex{1} \begin{tasks}(2)
  \task $\res{x=-2,\ x=- \sqrt{3},\ x=\sqrt{3},\ x=2}$
  \task $\res{x=-3,\ x=- \frac{2}{3},\ x=0,\ x=3,\ x=4}$
  \task $\res{x=-3,\ x=2}$
  \task $\res{x=-1,\ x=2,\ x=\frac{7}{2}}$
\end{tasks}


% =====================================================
\seccio{5}{Ecuaciones racionales}

\exf{2} \begin{tasks}(1)
  \task $8x^2+8x+8=7x^3\Rightarrow(x-2)(7x^2+6x+4)=0$ ($\Delta\lt0$): \ $\res{x=2}$
  \task Multiplicamos por $3(x-1)$: $2x^2+4x+9=11\Rightarrow x^2+2x-1=0$: \ $\res{x=-1\pm\sqrt2}$
  \task $2x(x+2)+3x(x-2)=6x^2\Rightarrow x^2+2x=0\Rightarrow x=0,\ x=-2$; \ $x=-2$ anula un denominador: \ $\res{x=0}$
  \task m.c.m. $=(x-1)(x-2)(x-4)$: $3(x-1)+(x-5)(x-2)=(x-1)(x-2)\Rightarrow x^2-4x+7=x^2-3x+2$: \ $\res{x=5}$
  \task Multiplicamos por $\sqrt2\,x$: $x^2+2=2x^2\Rightarrow$ $\res{x=\pm\sqrt2}$
  \task $-\dfrac{x}{x-2}=-x\Rightarrow x=x(x-2)\Rightarrow x^2-3x=0$: \ $\res{x=0,\ x=3}$
  \task Simplificando: $\dfrac{2x(x+1)}{(x-1)(x+3)}=2\Rightarrow x^2+x=x^2+2x-3$: \ $\res{x=3}$
  \task El primer miembro es $\dfrac{x+1}{x}$: es una \rest{identidad}, se cumple para todo $x\neq0,\ x\neq-1$.
\end{tasks}


% =====================================================
\seccio{6}{Ecuaciones con radicales}

\exf{2} \begin{tasks}(1)
  \task $2x-3=(1-x)^2\Rightarrow(x-2)^2=0$; \ $x=2$: $-1+1=0\neq2$: \ $\res{\text{no tiene solución}}$
  \task $2x^2+2x-3=(x+1)^2\Rightarrow x^2=4$; \ $x=2$ no vale ($-1-2\lt0$): \ $\res{x=-2}$
  \task $9(7x+1)=(5x-7)^2\Rightarrow25x^2-133x+40=0\Rightarrow x=5,\ x=\frac{8}{25}$; \ $x=\frac8{25}$ no vale: \ $\res{x=5}$
  \task $2x-1=12$: \ $\res{x=\frac{13}{2}}$
  \task $2-5x=3x^2\Rightarrow3x^2+5x-2=0\Rightarrow x=-2,\ x=\frac13$; \ $x=\frac13$ no vale: \ $\res{x=-2}$
  \task $2+\sqrt{x-4}=12-x\Rightarrow x-4=(10-x)^2\Rightarrow x^2-21x+104=0\Rightarrow x=8,\ x=13$; \ $x=13$ no vale: \ $\res{x=8}$
  \task $x=4-4\sqrt{x-4}+x-4\Rightarrow\sqrt{x-4}=0$: \ $\res{x=4}$
  \task $x-46=-14\sqrt{x+5}\Rightarrow x^2-288x+1136=0\Rightarrow x=4,\ x=284$; \ $x=284$ no vale: \ $\res{x=4}$
  \task $10\sqrt{x-1}=20\Rightarrow\sqrt{x-1}=2$: \ $\res{x=5}$
  \task $2\sqrt{2x^2-14x}=8-2x\Rightarrow x^2-6x-16=0\Rightarrow x=8,\ x=-2$; ninguna vale: \ $\res{\text{no tiene solución}}$
  \task $\sqrt{1-4x}=3\sqrt{3+x}\Rightarrow1-4x=27+9x$: \ $\res{x=-2}$
\end{tasks}

\ex{2} \begin{tasks}(1)
  \task $x^2-1=(x-1)^3\Rightarrow x(x-1)(x-3)=0$; todas valen ($\sqrt[3]{-1}=-1$): \ $\res{x=0,\ x=1,\ x=3}$
  \task $\sqrt{3x+7}=x+1\Rightarrow x^2-x-6=0\Rightarrow x=3,\ x=-2$; \ $x=-2$ no vale: \ $\res{x=3}$
  \task $4x-1=(x-4)^3\Rightarrow(x-7)(x^2-5x+9)=0$: \ $\res{x=7}$
  \task $(4-2x)^2=8x^2-16x\Rightarrow4x^2=16$; las dos valen: \ $\res{x=2,\ x=-2}$
  \task $(2x+2)^2=6x+10\Rightarrow2x^2+x-3=0\Rightarrow x=1,\ x=-\frac32$; \ $x=-\frac32$ no vale: \ $\res{x=1}$
\end{tasks}

\ex{2} $\sqrt{x+1}=2\sqrt{x-11}\Rightarrow x+1=4x-44$: \ $\res{x=15}$ \ (comprobación: $\sqrt{16}=4=2\sqrt4$).


% =====================================================
\seccio{7}{Ecuaciones logarítmicas}

\exf{2} \begin{tasks}(1)
  \task $2(x+3)=x+5$: \ $\res{x=-1}$
  \task $15-2x=x^2\Rightarrow x=3,\ x=-5$; \ $x=-5$ no vale: \ $\res{x=3}$
  \task $\dfrac{2x+3}{x-1}=36\Rightarrow34x=39$: \ $\res{x=\frac{39}{34}}$
  \task $x^2-3x+2=x^2\Rightarrow$ $\res{x=\frac23}$ (vale: $x\gt0$ y $x^2-3x+2=\frac49\gt0$)
  \task $\dfrac{2x-2}{x}=\dfrac{(x-1)^2}{x}\Rightarrow(x-1)(x-3)=0$; \ $x=1$ no vale: \ $\res{x=3}$
  \task $\log x^2=\log\dfrac{x}{20}\Rightarrow x^2=\dfrac x{20}$; \ $x=0$ no vale: \ $\res{x=\frac1{20}}$
  \task $2x(4-5x)=3(2x-x^2)\Rightarrow7x^2-2x=0$; \ $x=0$ no vale: \ $\res{x=\frac27}$
  \task $\dfrac{5(x+1)(x+2)}{x+2}=4\Rightarrow5x+5=4$: \ $\res{x=-\frac15}$
  \task $4-x=(x+2)^2\Rightarrow x^2+5x=0$; \ $x=-5$ no vale: \ $\res{x=0}$
  \task Cambio de base: $\dfrac{2\log3}{\log(x+3)}=\dfrac{\log3}{\log5}\Rightarrow\log(x+3)=2\log5\Rightarrow x+3=25$: \ $\res{x=22}$
  \task En base 3: $\log_332-\log_3(6-x)=2\log_3x\Rightarrow x^2(6-x)=32\Rightarrow(x+2)(x-4)^2=0$: \ $\res{x=4}$
\end{tasks}

\ex{1} $\log x+\log2=1\Rightarrow\log(2x)=\log10$: \ $\res{x=5}$

\ex{1} $2\log x=\log4+\log9\Rightarrow x^2=36$ ($x\gt0$): \ $\res{x=6}$


% =====================================================
\seccio{8}{Ecuaciones exponenciales}

\exf{2} \begin{tasks}(1)
  \task $2^{2x^2+2}=2^{5x+5}\Rightarrow2x^2-5x-3=0$: \ $\res{x=3,\ x=-\frac12}$
  \task $262144=4^9\Rightarrow(x-2)^2=9$: \ $\res{x=5,\ x=-1}$
  \task $x+4=\dfrac{\log5}{\log3}$: \ $\res{x=\frac{\log5}{\log3}-4\approx-2,535}$
  \task $2^{-x}=2^{2x(x-1)}\Rightarrow2x^2-x=0$: \ $\res{x=0,\ x=\frac12}$
  \task $13^{x^2+2x}=13^{-1}\Rightarrow(x+1)^2=0$: \ $\res{x=-1}$
  \task $2x^2-4x-2=4\Rightarrow x^2-2x-3=0$: \ $\res{x=3,\ x=-1}$
  \task $6^{4x-3}=12\Rightarrow4x-3=\dfrac{\log12}{\log6}$: \ $\res{x=\frac{3+\log_612}{4}\approx1,097}$
  \task $\left(\dfrac32\right)^x=\dfrac32$: \ $\res{x=1}$
  \task $3^{3/x}=3^{x+2}\Rightarrow x^2+2x-3=0$: \ $\res{x=1,\ x=-3}$
  \task $2^{x+4}=2^{3x}\Rightarrow x+4=3x$: \ $\res{x=2}$
  \task $x^2-5x+6=0$: \ $\res{x=2,\ x=3}$
  \task $7^{\frac{2(x-5)}{3}}=7^{-\frac{x-1}{2}}\Rightarrow4(x-5)=-3(x-1)$: \ $\res{x=\frac{23}{7}}$
\end{tasks}

\ex{2} \begin{tasks}(1)
  \task $2^x\left(2+1+\frac12\right)=28\Rightarrow2^x=8$: \ $\res{x=3}$
  \task $5^x\left(125-\frac15\right)=3120\Rightarrow5^x=25$: \ $\res{x=2}$
  \task $4^x\cdot\frac{21}{16}=336\Rightarrow4^x=256$: \ $\res{x=4}$
  \task $3^x\left(9+3+1+\frac13\right)=120\Rightarrow3^x=9$: \ $\res{x=2}$
  \task $2^x\cdot\frac{15}{16}=960\Rightarrow2^x=1024$: \ $\res{x=10}$
  \task $t=6^x$: $t^2+8t-9=0\Rightarrow t=1$ ($t=-9$ no vale): \ $\res{x=0}$
\end{tasks}

\ex{3} \begin{tasks}(1)
  \task $t=3^x$: $3t^2-10t+3=0\Rightarrow t=3,\ t=\frac13$: \ $\res{x=1,\ x=-1}$
  \task $t=10^{x+2}$: $2t^2+3t-5=0\Rightarrow t=1$ ($t=-\frac52$ no vale) $\Rightarrow x+2=0$: \ $\res{x=-2}$
  \task $t=3^x$: $t^2+5t-24=0\Rightarrow t=3$ ($t=-8$ no vale): \ $\res{x=1}$
  \task $t=3^x$: $t^2+\frac{t^2}{3}+\frac t3=111\Rightarrow4t^2+t-333=0\Rightarrow t=9$: \ $\res{x=2}$
  \task $t=3^x$: $9t+\frac{t^2}{9}=90\Rightarrow t^2+81t-810=0\Rightarrow t=9$ ($t=-90$ no vale): \ $\res{x=2}$
  \task $t=2^x$: $\frac{t^2}{16}-\frac{5t}{8}+1=0\Rightarrow t^2-10t+16=0\Rightarrow t=8,\ t=2$: \ $\res{x=3,\ x=1}$
  \task $t=3^x$: $81t^2+27t-810=0\Rightarrow3t^2+t-30=0\Rightarrow t=3$ ($t=-\frac{10}{3}$ no vale): \ $\res{x=1}$
  \task $t=2^x$: $4t^2+8t-320=0\Rightarrow t^2+2t-80=0\Rightarrow t=8$ ($t=-10$ no vale): \ $\res{x=3}$
\end{tasks}


% =====================================================
\seccio{9}{Problemas de ecuaciones}

\exf{2} $5x+2\cdot50+0\cdot50=4(x+100)\Rightarrow$ $\res{x=300\ \text{L}}$

\ex{1} $30\cdot5+10x=7(30+x)\Rightarrow3x=60$: \ $\res{20\ \text{L}}$

\ex{2} En 1 kg, $x$ kg de té blanco: $18x+14(1-x)=16,4\Rightarrow x=0,6$: \ $\res{60\,\%\ \text{blanco y }40\,\%\ \text{verde}}$

\ex{1} El 10\,\% son 15: \ $\res{150\ \text{alumnos}}$

\ex{2} A la venta queda el $75\,\%$ del $20\,\%$, es decir, el $15\,\%$: $0,15x=18\,000\Rightarrow$ $\res{120\,000\ \text{entradas}}$

\ex{1} África es el 15\,\%: $0,15x=3\Rightarrow$ $\res{20\ \text{millones de euros}}$

\ex{2} $5x-2(20-x)=58\Rightarrow7x=98$: \ $\res{14\ \text{aciertos y }6\ \text{errores}}$

\ex{3} $(x-20)\left(\dfrac{125}{x}+0,4\right)=147\Rightarrow x^2-75x-6250=0\Rightarrow x=125$ ($x=-50$ no tiene sentido): \ $\res{125\ \text{kg}}$

\ex{1} $x^2+(x+1)^2=25\Rightarrow x^2+x-12=0\Rightarrow x=3$: \ $\res{3\ \text{cm y }4\ \text{cm}}$

\ex{2} $x^2+(x+2)^2=(x+4)^2\Rightarrow x^2-4x-12=0\Rightarrow x=6$: \ $\res{6,\ 8\ \text{y }10\ \text{cm}}$

\ex{2} $x+y=70,\ x^2+y^2=2500\Rightarrow x^2-70x+1200=0$: \ $\res{40\ \text{m}\times30\ \text{m}}$

\ex{1} $\frac15+\frac1{20}=\frac14$: \ $\res{T=4\ \text{h}}$

\ex{1} $\frac12+\frac13=\frac56\Rightarrow T=\frac65$ h: \ $\res{1\ \text{h}\ 12\ \text{min}}$

\ex{2} $\dfrac1t+\dfrac1{2t}=\dfrac1{10}\Rightarrow t=15$: \ $\res{15\ \text{h y }30\ \text{h}}$

\ex{2} $\frac1{10}+\frac1{20}+\frac1{30}+\frac1{40}+\frac1{50}=\frac{137}{600}\Rightarrow T=\frac{600}{137}\approx\res{4,38\ \text{h}\approx4\ \text{h}\ 23\ \text{min}}$

\ex{3} $\dfrac{600}{x-10}=\dfrac{600}{x}+5\Rightarrow x^2-10x-1200=0\Rightarrow x=40$: \ $\res{40\ \text{personas}}$

\ex{2} \begin{tasks}(1)
  \task $\res{y(0)=300\ \text{enfermos}}$
  \task $200e^{0,2t}=800\Rightarrow t=5\ln4\approx6,93$: \ $\res{\text{el 7 de enero (hacia el final del día)}}$
\end{tasks}

\ex{2} \begin{tasks}(1)
  \task $e^{-0,1t}=\frac12\Rightarrow t=10\ln2\approx\res{6,93\ \text{días}}$
  \task $27e^{-1}\approx\res{9,93\ \text{mg}}$
\end{tasks}

\ex{2} $2^t=\dfrac{100}{9}\Rightarrow t=\dfrac{\log(100/9)}{\log2}\approx\res{3,47\ \text{h}\approx3\ \text{h}\ 28\ \text{min}}$

\ex{3} \begin{tasks}(1)
  \task $1+Ka=4$, \ $1+Ka^2=2,5\Rightarrow Ka=3,\ Ka^2=1,5$: \ $\res{a=\frac12,\ K=6}$
  \task $1+6\cdot0,5^{x/10}=\frac43\Rightarrow0,5^{x/10}=\frac1{18}\Rightarrow x=10\log_218\approx\res{41,7\ \text{años}}$
\end{tasks}


% =====================================================
\seccio{10}{Inecuaciones con una incógnita}

\exf{1} \begin{tasks}(1)
  \task $-8+4-2=-6\le6$: \rest{sí}
  \task $\frac1e-1\approx-0,63\not\ge0$: \rest{no}
  \task $2^{-0,5}-2,5\approx-1,79\lt3^{-0,5}\approx0,58$: \rest{sí}
\end{tasks}

\ex{1} \begin{tasks}(1)
  \task $4x-12-x+2\le4x\Rightarrow-x\le10$: \ $\res{x\in [-10,\, +\infty)}$\\[2pt]\begin{tikzpicture}[x=0.420cm,baseline=-3pt]
\draw[-latex,muted] (-13.6,0)--(-6.3,0);
\draw[muted] (-13,0.1)--(-13,-0.1) node[below=-1pt,font=\scriptsize,text=muted]{$-13$};
\draw[muted] (-12,0.1)--(-12,-0.1) node[below=-1pt,font=\scriptsize,text=muted]{$-12$};
\draw[muted] (-11,0.1)--(-11,-0.1) node[below=-1pt,font=\scriptsize,text=muted]{$-11$};
\draw[muted] (-10,0.1)--(-10,-0.1) node[below=-1pt,font=\scriptsize,text=muted]{$-10$};
\draw[muted] (-9,0.1)--(-9,-0.1) node[below=-1pt,font=\scriptsize,text=muted]{$-9$};
\draw[muted] (-8,0.1)--(-8,-0.1) node[below=-1pt,font=\scriptsize,text=muted]{$-8$};
\draw[muted] (-7,0.1)--(-7,-0.1) node[below=-1pt,font=\scriptsize,text=muted]{$-7$};
\draw[accent,line width=2.6pt] (-10,0)--(-6.55,0);
\fill[accent] (-6.25,0)--(-6.65,0.17)--(-6.65,-0.17)--cycle;
\fill[accent] (-10,0) circle (2.3pt);
\end{tikzpicture}
  \task $9x-18\gt9x+1\Rightarrow-18\gt1$, falso: \ $\res{\text{no tiene solución}}$
  \task $6x+8\lt\frac{x+38}{2}\Rightarrow11x\lt22$: \ $\res{x\in (-\infty,\, 2)}$\\[2pt]\begin{tikzpicture}[x=0.420cm,baseline=-3pt]
\draw[-latex,muted] (-1.6,0)--(5.7,0);
\draw[muted] (-1,0.1)--(-1,-0.1) node[below=-1pt,font=\scriptsize,text=muted]{$-1$};
\draw[muted] (0,0.1)--(0,-0.1) node[below=-1pt,font=\scriptsize,text=muted]{$0$};
\draw[muted] (1,0.1)--(1,-0.1) node[below=-1pt,font=\scriptsize,text=muted]{$1$};
\draw[muted] (2,0.1)--(2,-0.1) node[below=-1pt,font=\scriptsize,text=muted]{$2$};
\draw[muted] (3,0.1)--(3,-0.1) node[below=-1pt,font=\scriptsize,text=muted]{$3$};
\draw[muted] (4,0.1)--(4,-0.1) node[below=-1pt,font=\scriptsize,text=muted]{$4$};
\draw[muted] (5,0.1)--(5,-0.1) node[below=-1pt,font=\scriptsize,text=muted]{$5$};
\draw[accent,line width=2.6pt] (-1.45,0)--(2,0);
\fill[accent] (-1.75,0)--(-1.35,0.17)--(-1.35,-0.17)--cycle;
\filldraw[fill=white,draw=accent,line width=1.1pt] (2,0) circle (2.3pt);
\end{tikzpicture}
  \task $2x^2-8x\le0\Rightarrow2x(x-4)\le0$: \ $\res{x\in [0,\, 4]}$\\[2pt]\begin{tikzpicture}[x=0.420cm,baseline=-3pt]
\draw[-latex,muted] (-1.6,0)--(5.7,0);
\draw[muted] (-1,0.1)--(-1,-0.1) node[below=-1pt,font=\scriptsize,text=muted]{$-1$};
\draw[muted] (0,0.1)--(0,-0.1) node[below=-1pt,font=\scriptsize,text=muted]{$0$};
\draw[muted] (1,0.1)--(1,-0.1) node[below=-1pt,font=\scriptsize,text=muted]{$1$};
\draw[muted] (2,0.1)--(2,-0.1) node[below=-1pt,font=\scriptsize,text=muted]{$2$};
\draw[muted] (3,0.1)--(3,-0.1) node[below=-1pt,font=\scriptsize,text=muted]{$3$};
\draw[muted] (4,0.1)--(4,-0.1) node[below=-1pt,font=\scriptsize,text=muted]{$4$};
\draw[muted] (5,0.1)--(5,-0.1) node[below=-1pt,font=\scriptsize,text=muted]{$5$};
\draw[accent,line width=2.6pt] (0,0)--(4,0);
\fill[accent] (0,0) circle (2.3pt);
\fill[accent] (4,0) circle (2.3pt);
\end{tikzpicture}
\end{tasks}

\ex{2} \begin{tasks}(1)
  \task $- x (2 x^{2} - 3)\gt0$: \ $\res{x\in (-\infty,\, {-} \frac{\sqrt{6}}{2})\cup(0,\, \frac{\sqrt{6}}{2})}$\\[2pt]\begin{tikzpicture}[x=0.420cm,baseline=-3pt]
\draw[-latex,muted] (-3.6,0)--(3.7,0);
\draw[muted] (-3,0.1)--(-3,-0.1) node[below=-1pt,font=\scriptsize,text=muted]{$-3$};
\draw[muted] (-2,0.1)--(-2,-0.1) node[below=-1pt,font=\scriptsize,text=muted]{$-2$};
\draw[muted] (-1,0.1)--(-1,-0.1) node[below=-1pt,font=\scriptsize,text=muted]{$-1$};
\draw[muted] (0,0.1)--(0,-0.1) node[below=-1pt,font=\scriptsize,text=muted]{$0$};
\draw[muted] (1,0.1)--(1,-0.1) node[below=-1pt,font=\scriptsize,text=muted]{$1$};
\draw[muted] (2,0.1)--(2,-0.1) node[below=-1pt,font=\scriptsize,text=muted]{$2$};
\draw[muted] (3,0.1)--(3,-0.1) node[below=-1pt,font=\scriptsize,text=muted]{$3$};
\draw[accent,line width=2.6pt] (-3.45,0)--(-1.22474,0);
\fill[accent] (-3.75,0)--(-3.35,0.17)--(-3.35,-0.17)--cycle;
\filldraw[fill=white,draw=accent,line width=1.1pt] (-1.22474,0) circle (2.3pt);
\node[above=1pt,font=\scriptsize,text=accent] at (-1.22474,0) {$- \frac{\sqrt{6}}{2}$};
\draw[accent,line width=2.6pt] (0,0)--(1.22474,0);
\filldraw[fill=white,draw=accent,line width=1.1pt] (0,0) circle (2.3pt);
\filldraw[fill=white,draw=accent,line width=1.1pt] (1.22474,0) circle (2.3pt);
\node[above=1pt,font=\scriptsize,text=accent] at (1.22474,0) {$\frac{\sqrt{6}}{2}$};
\end{tikzpicture}
  \task $(x + 3) (3 x - 8)\ge0$: \ $\res{x\in (-\infty,\, {-}3]\cup[\frac{8}{3},\, +\infty)}$\\[2pt]\begin{tikzpicture}[x=0.420cm,baseline=-3pt]
\draw[-latex,muted] (-4.6,0)--(4.7,0);
\draw[muted] (-4,0.1)--(-4,-0.1) node[below=-1pt,font=\scriptsize,text=muted]{$-4$};
\draw[muted] (-3,0.1)--(-3,-0.1) node[below=-1pt,font=\scriptsize,text=muted]{$-3$};
\draw[muted] (-2,0.1)--(-2,-0.1) node[below=-1pt,font=\scriptsize,text=muted]{$-2$};
\draw[muted] (-1,0.1)--(-1,-0.1) node[below=-1pt,font=\scriptsize,text=muted]{$-1$};
\draw[muted] (0,0.1)--(0,-0.1) node[below=-1pt,font=\scriptsize,text=muted]{$0$};
\draw[muted] (1,0.1)--(1,-0.1) node[below=-1pt,font=\scriptsize,text=muted]{$1$};
\draw[muted] (2,0.1)--(2,-0.1) node[below=-1pt,font=\scriptsize,text=muted]{$2$};
\draw[muted] (3,0.1)--(3,-0.1) node[below=-1pt,font=\scriptsize,text=muted]{$3$};
\draw[muted] (4,0.1)--(4,-0.1) node[below=-1pt,font=\scriptsize,text=muted]{$4$};
\draw[accent,line width=2.6pt] (-4.45,0)--(-3,0);
\fill[accent] (-4.75,0)--(-4.35,0.17)--(-4.35,-0.17)--cycle;
\fill[accent] (-3,0) circle (2.3pt);
\draw[accent,line width=2.6pt] (2.66667,0)--(4.45,0);
\fill[accent] (4.75,0)--(4.35,0.17)--(4.35,-0.17)--cycle;
\fill[accent] (2.66667,0) circle (2.3pt);
\node[above=1pt,font=\scriptsize,text=accent] at (2.66667,0) {$\frac{8}{3}$};
\end{tikzpicture}
  \task $x (x - 2) (x + 2)\gt0$: \ $\res{x\in (-2,\, 0)\cup(2,\, +\infty)}$\\[2pt]\begin{tikzpicture}[x=0.420cm,baseline=-3pt]
\draw[-latex,muted] (-3.6,0)--(3.7,0);
\draw[muted] (-3,0.1)--(-3,-0.1) node[below=-1pt,font=\scriptsize,text=muted]{$-3$};
\draw[muted] (-2,0.1)--(-2,-0.1) node[below=-1pt,font=\scriptsize,text=muted]{$-2$};
\draw[muted] (-1,0.1)--(-1,-0.1) node[below=-1pt,font=\scriptsize,text=muted]{$-1$};
\draw[muted] (0,0.1)--(0,-0.1) node[below=-1pt,font=\scriptsize,text=muted]{$0$};
\draw[muted] (1,0.1)--(1,-0.1) node[below=-1pt,font=\scriptsize,text=muted]{$1$};
\draw[muted] (2,0.1)--(2,-0.1) node[below=-1pt,font=\scriptsize,text=muted]{$2$};
\draw[muted] (3,0.1)--(3,-0.1) node[below=-1pt,font=\scriptsize,text=muted]{$3$};
\draw[accent,line width=2.6pt] (-2,0)--(0,0);
\filldraw[fill=white,draw=accent,line width=1.1pt] (-2,0) circle (2.3pt);
\filldraw[fill=white,draw=accent,line width=1.1pt] (0,0) circle (2.3pt);
\draw[accent,line width=2.6pt] (2,0)--(3.45,0);
\fill[accent] (3.75,0)--(3.35,0.17)--(3.35,-0.17)--cycle;
\filldraw[fill=white,draw=accent,line width=1.1pt] (2,0) circle (2.3pt);
\end{tikzpicture}
  \task $(x - 1) (x + 1)^{2}\le0$: \ $\res{x\in (-\infty,\, 1]}$\\[2pt]\begin{tikzpicture}[x=0.420cm,baseline=-3pt]
\draw[-latex,muted] (-2.6,0)--(4.7,0);
\draw[muted] (-2,0.1)--(-2,-0.1) node[below=-1pt,font=\scriptsize,text=muted]{$-2$};
\draw[muted] (-1,0.1)--(-1,-0.1) node[below=-1pt,font=\scriptsize,text=muted]{$-1$};
\draw[muted] (0,0.1)--(0,-0.1) node[below=-1pt,font=\scriptsize,text=muted]{$0$};
\draw[muted] (1,0.1)--(1,-0.1) node[below=-1pt,font=\scriptsize,text=muted]{$1$};
\draw[muted] (2,0.1)--(2,-0.1) node[below=-1pt,font=\scriptsize,text=muted]{$2$};
\draw[muted] (3,0.1)--(3,-0.1) node[below=-1pt,font=\scriptsize,text=muted]{$3$};
\draw[muted] (4,0.1)--(4,-0.1) node[below=-1pt,font=\scriptsize,text=muted]{$4$};
\draw[accent,line width=2.6pt] (-2.45,0)--(1,0);
\fill[accent] (-2.75,0)--(-2.35,0.17)--(-2.35,-0.17)--cycle;
\fill[accent] (1,0) circle (2.3pt);
\end{tikzpicture}
  \task $(x - 1) (7 x - 1)\ge0$: \ $\res{x\in (-\infty,\, \frac{1}{7}]\cup[1,\, +\infty)}$\\[2pt]\begin{tikzpicture}[x=0.420cm,baseline=-3pt]
\draw[-latex,muted] (-3.6,0)--(4.7,0);
\draw[muted] (-3,0.1)--(-3,-0.1) node[below=-1pt,font=\scriptsize,text=muted]{$-3$};
\draw[muted] (-2,0.1)--(-2,-0.1) node[below=-1pt,font=\scriptsize,text=muted]{$-2$};
\draw[muted] (-1,0.1)--(-1,-0.1) node[below=-1pt,font=\scriptsize,text=muted]{$-1$};
\draw[muted] (0,0.1)--(0,-0.1) node[below=-1pt,font=\scriptsize,text=muted]{$0$};
\draw[muted] (1,0.1)--(1,-0.1) node[below=-1pt,font=\scriptsize,text=muted]{$1$};
\draw[muted] (2,0.1)--(2,-0.1) node[below=-1pt,font=\scriptsize,text=muted]{$2$};
\draw[muted] (3,0.1)--(3,-0.1) node[below=-1pt,font=\scriptsize,text=muted]{$3$};
\draw[muted] (4,0.1)--(4,-0.1) node[below=-1pt,font=\scriptsize,text=muted]{$4$};
\draw[accent,line width=2.6pt] (-3.45,0)--(0.142857,0);
\fill[accent] (-3.75,0)--(-3.35,0.17)--(-3.35,-0.17)--cycle;
\fill[accent] (0.142857,0) circle (2.3pt);
\node[above=1pt,font=\scriptsize,text=accent] at (0.142857,0) {$\frac{1}{7}$};
\draw[accent,line width=2.6pt] (1,0)--(4.45,0);
\fill[accent] (4.75,0)--(4.35,0.17)--(4.35,-0.17)--cycle;
\fill[accent] (1,0) circle (2.3pt);
\end{tikzpicture}
  \task $- (x - 1) (x + 1) (x + 2)\le0$: \ $\res{x\in [-2,\, {-}1]\cup[1,\, +\infty)}$\\[2pt]\begin{tikzpicture}[x=0.420cm,baseline=-3pt]
\draw[-latex,muted] (-4.6,0)--(3.7,0);
\draw[muted] (-4,0.1)--(-4,-0.1) node[below=-1pt,font=\scriptsize,text=muted]{$-4$};
\draw[muted] (-3,0.1)--(-3,-0.1) node[below=-1pt,font=\scriptsize,text=muted]{$-3$};
\draw[muted] (-2,0.1)--(-2,-0.1) node[below=-1pt,font=\scriptsize,text=muted]{$-2$};
\draw[muted] (-1,0.1)--(-1,-0.1) node[below=-1pt,font=\scriptsize,text=muted]{$-1$};
\draw[muted] (0,0.1)--(0,-0.1) node[below=-1pt,font=\scriptsize,text=muted]{$0$};
\draw[muted] (1,0.1)--(1,-0.1) node[below=-1pt,font=\scriptsize,text=muted]{$1$};
\draw[muted] (2,0.1)--(2,-0.1) node[below=-1pt,font=\scriptsize,text=muted]{$2$};
\draw[muted] (3,0.1)--(3,-0.1) node[below=-1pt,font=\scriptsize,text=muted]{$3$};
\draw[accent,line width=2.6pt] (-2,0)--(-1,0);
\fill[accent] (-2,0) circle (2.3pt);
\fill[accent] (-1,0) circle (2.3pt);
\draw[accent,line width=2.6pt] (1,0)--(3.45,0);
\fill[accent] (3.75,0)--(3.35,0.17)--(3.35,-0.17)--cycle;
\fill[accent] (1,0) circle (2.3pt);
\end{tikzpicture}
  \task $(x - 1) (x + 1) (x^{2} + 1)\gt0$: \ $\res{x\in (-\infty,\, {-}1)\cup(1,\, +\infty)}$\\[2pt]\begin{tikzpicture}[x=0.420cm,baseline=-3pt]
\draw[-latex,muted] (-3.6,0)--(3.7,0);
\draw[muted] (-3,0.1)--(-3,-0.1) node[below=-1pt,font=\scriptsize,text=muted]{$-3$};
\draw[muted] (-2,0.1)--(-2,-0.1) node[below=-1pt,font=\scriptsize,text=muted]{$-2$};
\draw[muted] (-1,0.1)--(-1,-0.1) node[below=-1pt,font=\scriptsize,text=muted]{$-1$};
\draw[muted] (0,0.1)--(0,-0.1) node[below=-1pt,font=\scriptsize,text=muted]{$0$};
\draw[muted] (1,0.1)--(1,-0.1) node[below=-1pt,font=\scriptsize,text=muted]{$1$};
\draw[muted] (2,0.1)--(2,-0.1) node[below=-1pt,font=\scriptsize,text=muted]{$2$};
\draw[muted] (3,0.1)--(3,-0.1) node[below=-1pt,font=\scriptsize,text=muted]{$3$};
\draw[accent,line width=2.6pt] (-3.45,0)--(-1,0);
\fill[accent] (-3.75,0)--(-3.35,0.17)--(-3.35,-0.17)--cycle;
\filldraw[fill=white,draw=accent,line width=1.1pt] (-1,0) circle (2.3pt);
\draw[accent,line width=2.6pt] (1,0)--(3.45,0);
\fill[accent] (3.75,0)--(3.35,0.17)--(3.35,-0.17)--cycle;
\filldraw[fill=white,draw=accent,line width=1.1pt] (1,0) circle (2.3pt);
\end{tikzpicture}
  \task $(x + 1) (x + 2)\lt0$: \ $\res{x\in (-2,\, {-}1)}$\\[2pt]\begin{tikzpicture}[x=0.420cm,baseline=-3pt]
\draw[-latex,muted] (-5.6,0)--(2.7,0);
\draw[muted] (-5,0.1)--(-5,-0.1) node[below=-1pt,font=\scriptsize,text=muted]{$-5$};
\draw[muted] (-4,0.1)--(-4,-0.1) node[below=-1pt,font=\scriptsize,text=muted]{$-4$};
\draw[muted] (-3,0.1)--(-3,-0.1) node[below=-1pt,font=\scriptsize,text=muted]{$-3$};
\draw[muted] (-2,0.1)--(-2,-0.1) node[below=-1pt,font=\scriptsize,text=muted]{$-2$};
\draw[muted] (-1,0.1)--(-1,-0.1) node[below=-1pt,font=\scriptsize,text=muted]{$-1$};
\draw[muted] (0,0.1)--(0,-0.1) node[below=-1pt,font=\scriptsize,text=muted]{$0$};
\draw[muted] (1,0.1)--(1,-0.1) node[below=-1pt,font=\scriptsize,text=muted]{$1$};
\draw[muted] (2,0.1)--(2,-0.1) node[below=-1pt,font=\scriptsize,text=muted]{$2$};
\draw[accent,line width=2.6pt] (-2,0)--(-1,0);
\filldraw[fill=white,draw=accent,line width=1.1pt] (-2,0) circle (2.3pt);
\filldraw[fill=white,draw=accent,line width=1.1pt] (-1,0) circle (2.3pt);
\end{tikzpicture}
\end{tasks}

\ex{2} \begin{tasks}(1)
  \task $\dfrac{5 x - 2}{2 x + 1}\ge0$: \ $\res{x\in (-\infty,\, {-} \frac{1}{2})\cup[\frac{2}{5},\, +\infty)}$\\[2pt]\begin{tikzpicture}[x=0.420cm,baseline=-3pt]
\draw[-latex,muted] (-3.6,0)--(3.7,0);
\draw[muted] (-3,0.1)--(-3,-0.1) node[below=-1pt,font=\scriptsize,text=muted]{$-3$};
\draw[muted] (-2,0.1)--(-2,-0.1) node[below=-1pt,font=\scriptsize,text=muted]{$-2$};
\draw[muted] (-1,0.1)--(-1,-0.1) node[below=-1pt,font=\scriptsize,text=muted]{$-1$};
\draw[muted] (0,0.1)--(0,-0.1) node[below=-1pt,font=\scriptsize,text=muted]{$0$};
\draw[muted] (1,0.1)--(1,-0.1) node[below=-1pt,font=\scriptsize,text=muted]{$1$};
\draw[muted] (2,0.1)--(2,-0.1) node[below=-1pt,font=\scriptsize,text=muted]{$2$};
\draw[muted] (3,0.1)--(3,-0.1) node[below=-1pt,font=\scriptsize,text=muted]{$3$};
\draw[accent,line width=2.6pt] (-3.45,0)--(-0.5,0);
\fill[accent] (-3.75,0)--(-3.35,0.17)--(-3.35,-0.17)--cycle;
\filldraw[fill=white,draw=accent,line width=1.1pt] (-0.5,0) circle (2.3pt);
\node[above=1pt,font=\scriptsize,text=accent] at (-0.5,0) {$- \frac{1}{2}$};
\draw[accent,line width=2.6pt] (0.4,0)--(3.45,0);
\fill[accent] (3.75,0)--(3.35,0.17)--(3.35,-0.17)--cycle;
\fill[accent] (0.4,0) circle (2.3pt);
\node[above=1pt,font=\scriptsize,text=accent] at (0.4,0) {$\frac{2}{5}$};
\end{tikzpicture}
  \task $\dfrac{(x - 2) (x + 1)}{(x - 1) (x + 2)}\ge0$: \ $\res{x\in (-\infty,\, {-}2)\cup[-1,\, 1)\cup[2,\, +\infty)}$\\[2pt]\begin{tikzpicture}[x=0.420cm,baseline=-3pt]
\draw[-latex,muted] (-3.6,0)--(3.7,0);
\draw[muted] (-3,0.1)--(-3,-0.1) node[below=-1pt,font=\scriptsize,text=muted]{$-3$};
\draw[muted] (-2,0.1)--(-2,-0.1) node[below=-1pt,font=\scriptsize,text=muted]{$-2$};
\draw[muted] (-1,0.1)--(-1,-0.1) node[below=-1pt,font=\scriptsize,text=muted]{$-1$};
\draw[muted] (0,0.1)--(0,-0.1) node[below=-1pt,font=\scriptsize,text=muted]{$0$};
\draw[muted] (1,0.1)--(1,-0.1) node[below=-1pt,font=\scriptsize,text=muted]{$1$};
\draw[muted] (2,0.1)--(2,-0.1) node[below=-1pt,font=\scriptsize,text=muted]{$2$};
\draw[muted] (3,0.1)--(3,-0.1) node[below=-1pt,font=\scriptsize,text=muted]{$3$};
\draw[accent,line width=2.6pt] (-3.45,0)--(-2,0);
\fill[accent] (-3.75,0)--(-3.35,0.17)--(-3.35,-0.17)--cycle;
\filldraw[fill=white,draw=accent,line width=1.1pt] (-2,0) circle (2.3pt);
\draw[accent,line width=2.6pt] (-1,0)--(1,0);
\fill[accent] (-1,0) circle (2.3pt);
\filldraw[fill=white,draw=accent,line width=1.1pt] (1,0) circle (2.3pt);
\draw[accent,line width=2.6pt] (2,0)--(3.45,0);
\fill[accent] (3.75,0)--(3.35,0.17)--(3.35,-0.17)--cycle;
\fill[accent] (2,0) circle (2.3pt);
\end{tikzpicture}
  \task $\dfrac{4 x - 5}{(x + 1) (4 x - 5)}\lt0$: \ $\res{x\in (-\infty,\, {-}1)}$\\[2pt]\begin{tikzpicture}[x=0.420cm,baseline=-3pt]
\draw[-latex,muted] (-4.6,0)--(2.7,0);
\draw[muted] (-4,0.1)--(-4,-0.1) node[below=-1pt,font=\scriptsize,text=muted]{$-4$};
\draw[muted] (-3,0.1)--(-3,-0.1) node[below=-1pt,font=\scriptsize,text=muted]{$-3$};
\draw[muted] (-2,0.1)--(-2,-0.1) node[below=-1pt,font=\scriptsize,text=muted]{$-2$};
\draw[muted] (-1,0.1)--(-1,-0.1) node[below=-1pt,font=\scriptsize,text=muted]{$-1$};
\draw[muted] (0,0.1)--(0,-0.1) node[below=-1pt,font=\scriptsize,text=muted]{$0$};
\draw[muted] (1,0.1)--(1,-0.1) node[below=-1pt,font=\scriptsize,text=muted]{$1$};
\draw[muted] (2,0.1)--(2,-0.1) node[below=-1pt,font=\scriptsize,text=muted]{$2$};
\draw[accent,line width=2.6pt] (-4.45,0)--(-1,0);
\fill[accent] (-4.75,0)--(-4.35,0.17)--(-4.35,-0.17)--cycle;
\filldraw[fill=white,draw=accent,line width=1.1pt] (-1,0) circle (2.3pt);
\end{tikzpicture}
  \task $\dfrac{x^{3}}{x - 9}\ge0$: \ $\res{x\in (-\infty,\, 0]\cup(9,\, +\infty)}$\\[2pt]\begin{tikzpicture}[x=0.420cm,baseline=-3pt]
\draw[-latex,muted] (-1.6,0)--(10.7,0);
\draw[muted] (-1,0.1)--(-1,-0.1) node[below=-1pt,font=\scriptsize,text=muted]{$-1$};
\draw[muted] (0,0.1)--(0,-0.1) node[below=-1pt,font=\scriptsize,text=muted]{$0$};
\draw[muted] (1,0.1)--(1,-0.1) node[below=-1pt,font=\scriptsize,text=muted]{$1$};
\draw[muted] (2,0.1)--(2,-0.1) node[below=-1pt,font=\scriptsize,text=muted]{$2$};
\draw[muted] (3,0.1)--(3,-0.1) node[below=-1pt,font=\scriptsize,text=muted]{$3$};
\draw[muted] (4,0.1)--(4,-0.1) node[below=-1pt,font=\scriptsize,text=muted]{$4$};
\draw[muted] (5,0.1)--(5,-0.1) node[below=-1pt,font=\scriptsize,text=muted]{$5$};
\draw[muted] (6,0.1)--(6,-0.1) node[below=-1pt,font=\scriptsize,text=muted]{$6$};
\draw[muted] (7,0.1)--(7,-0.1) node[below=-1pt,font=\scriptsize,text=muted]{$7$};
\draw[muted] (8,0.1)--(8,-0.1) node[below=-1pt,font=\scriptsize,text=muted]{$8$};
\draw[muted] (9,0.1)--(9,-0.1) node[below=-1pt,font=\scriptsize,text=muted]{$9$};
\draw[muted] (10,0.1)--(10,-0.1) node[below=-1pt,font=\scriptsize,text=muted]{$10$};
\draw[accent,line width=2.6pt] (-1.45,0)--(0,0);
\fill[accent] (-1.75,0)--(-1.35,0.17)--(-1.35,-0.17)--cycle;
\fill[accent] (0,0) circle (2.3pt);
\draw[accent,line width=2.6pt] (9,0)--(10.45,0);
\fill[accent] (10.75,0)--(10.35,0.17)--(10.35,-0.17)--cycle;
\filldraw[fill=white,draw=accent,line width=1.1pt] (9,0) circle (2.3pt);
\end{tikzpicture}
  \task $\dfrac{(x - 1) (x + 5)}{x (x - 1) (x + 1)}\ge0$: \ $\res{x\in [-5,\, {-}1)\cup(0,\, 1)\cup(1,\, +\infty)}$\\[2pt]\begin{tikzpicture}[x=0.420cm,baseline=-3pt]
\draw[-latex,muted] (-6.6,0)--(2.7,0);
\draw[muted] (-6,0.1)--(-6,-0.1) node[below=-1pt,font=\scriptsize,text=muted]{$-6$};
\draw[muted] (-5,0.1)--(-5,-0.1) node[below=-1pt,font=\scriptsize,text=muted]{$-5$};
\draw[muted] (-4,0.1)--(-4,-0.1) node[below=-1pt,font=\scriptsize,text=muted]{$-4$};
\draw[muted] (-3,0.1)--(-3,-0.1) node[below=-1pt,font=\scriptsize,text=muted]{$-3$};
\draw[muted] (-2,0.1)--(-2,-0.1) node[below=-1pt,font=\scriptsize,text=muted]{$-2$};
\draw[muted] (-1,0.1)--(-1,-0.1) node[below=-1pt,font=\scriptsize,text=muted]{$-1$};
\draw[muted] (0,0.1)--(0,-0.1) node[below=-1pt,font=\scriptsize,text=muted]{$0$};
\draw[muted] (1,0.1)--(1,-0.1) node[below=-1pt,font=\scriptsize,text=muted]{$1$};
\draw[muted] (2,0.1)--(2,-0.1) node[below=-1pt,font=\scriptsize,text=muted]{$2$};
\draw[accent,line width=2.6pt] (-5,0)--(-1,0);
\fill[accent] (-5,0) circle (2.3pt);
\filldraw[fill=white,draw=accent,line width=1.1pt] (-1,0) circle (2.3pt);
\draw[accent,line width=2.6pt] (0,0)--(1,0);
\filldraw[fill=white,draw=accent,line width=1.1pt] (0,0) circle (2.3pt);
\filldraw[fill=white,draw=accent,line width=1.1pt] (1,0) circle (2.3pt);
\draw[accent,line width=2.6pt] (1,0)--(2.45,0);
\fill[accent] (2.75,0)--(2.35,0.17)--(2.35,-0.17)--cycle;
\filldraw[fill=white,draw=accent,line width=1.1pt] (1,0) circle (2.3pt);
\end{tikzpicture}
  \task $\dfrac{2 x + 3}{(x + 1) (x^{2} - x + 1)}\ge0$: \ $\res{x\in (-\infty,\, {-} \frac{3}{2}]\cup(-1,\, +\infty)}$\\[2pt]\begin{tikzpicture}[x=0.420cm,baseline=-3pt]
\draw[-latex,muted] (-5.6,0)--(2.7,0);
\draw[muted] (-5,0.1)--(-5,-0.1) node[below=-1pt,font=\scriptsize,text=muted]{$-5$};
\draw[muted] (-4,0.1)--(-4,-0.1) node[below=-1pt,font=\scriptsize,text=muted]{$-4$};
\draw[muted] (-3,0.1)--(-3,-0.1) node[below=-1pt,font=\scriptsize,text=muted]{$-3$};
\draw[muted] (-2,0.1)--(-2,-0.1) node[below=-1pt,font=\scriptsize,text=muted]{$-2$};
\draw[muted] (-1,0.1)--(-1,-0.1) node[below=-1pt,font=\scriptsize,text=muted]{$-1$};
\draw[muted] (0,0.1)--(0,-0.1) node[below=-1pt,font=\scriptsize,text=muted]{$0$};
\draw[muted] (1,0.1)--(1,-0.1) node[below=-1pt,font=\scriptsize,text=muted]{$1$};
\draw[muted] (2,0.1)--(2,-0.1) node[below=-1pt,font=\scriptsize,text=muted]{$2$};
\draw[accent,line width=2.6pt] (-5.45,0)--(-1.5,0);
\fill[accent] (-5.75,0)--(-5.35,0.17)--(-5.35,-0.17)--cycle;
\fill[accent] (-1.5,0) circle (2.3pt);
\node[above=1pt,font=\scriptsize,text=accent] at (-1.5,0) {$- \frac{3}{2}$};
\draw[accent,line width=2.6pt] (-1,0)--(2.45,0);
\fill[accent] (2.75,0)--(2.35,0.17)--(2.35,-0.17)--cycle;
\filldraw[fill=white,draw=accent,line width=1.1pt] (-1,0) circle (2.3pt);
\end{tikzpicture}
\end{tasks}


% =====================================================
\seccio{11}{Problemas de inecuaciones}

\exf{1} $10+0,05x\lt0,12x\Rightarrow x\gt\frac{1000}{7}\approx142,9$: \ $\res{\text{a partir de 143 minutos}}$

\ex{1} $1,2+0,7x\ge5\Rightarrow x\ge\frac{38}{7}\approx5,43$: \ $\res{x\in [\frac{38}{7},\, 10]}$\\[2pt]\begin{tikzpicture}[x=0.420cm,baseline=-3pt]
\draw[-latex,muted] (3.4,0)--(11.7,0);
\draw[muted] (4,0.1)--(4,-0.1) node[below=-1pt,font=\scriptsize,text=muted]{$4$};
\draw[muted] (5,0.1)--(5,-0.1) node[below=-1pt,font=\scriptsize,text=muted]{$5$};
\draw[muted] (6,0.1)--(6,-0.1) node[below=-1pt,font=\scriptsize,text=muted]{$6$};
\draw[muted] (7,0.1)--(7,-0.1) node[below=-1pt,font=\scriptsize,text=muted]{$7$};
\draw[muted] (8,0.1)--(8,-0.1) node[below=-1pt,font=\scriptsize,text=muted]{$8$};
\draw[muted] (9,0.1)--(9,-0.1) node[below=-1pt,font=\scriptsize,text=muted]{$9$};
\draw[muted] (10,0.1)--(10,-0.1) node[below=-1pt,font=\scriptsize,text=muted]{$10$};
\draw[muted] (11,0.1)--(11,-0.1) node[below=-1pt,font=\scriptsize,text=muted]{$11$};
\draw[accent,line width=2.6pt] (5.42857,0)--(10,0);
\fill[accent] (5.42857,0) circle (2.3pt);
\node[above=1pt,font=\scriptsize,text=accent] at (5.42857,0) {$\frac{38}{7}$};
\fill[accent] (10,0) circle (2.3pt);
\end{tikzpicture}

\ex{2} $20t-5t^2\gt15\Rightarrow t^2-4t+3\lt0\Rightarrow t\in(1,3)$: \ $\res{2\ \text{segundos}}$


% =====================================================
\seccio{12}{Ampliación y problemas}

\exf{1} $P(-1)=1+2+1=$ $\res{4}$

\ex{2} $P(x)=a(x-1)(x+2)(x-3)$; \ $P(0)=6a=12\Rightarrow a=2$: \ $\res{2 x^{3} - 4 x^{2} - 10 x + 12}$

\ex{2} \begin{tasks}(2)
  \task $\dfrac{(x-3)(x^2+3x+9)}{(x-3)(x+3)}=\res{\frac{x^2+3x+9}{x+3}}$
  \task $\dfrac{\frac{x-1}{x}}{\frac{(x-1)(x+1)}{x^2}}=\res{\frac{x}{x+1}}$
\end{tasks}

\ex{2} $2x^2+2(x+1)^2=5x(x+1)\Rightarrow x^2+x-2=0$: \ $\res{x=1,\ x=-2}$ (ninguna anula los denominadores).

\ex{2} $x+3=1+2\sqrt{x-2}+x-2\Rightarrow\sqrt{x-2}=2$: \ $\res{x=6}$ \ ($3-2=1$ \checkmark).

\ex{2} $t=2^x$: $2t^2-5t+2=0\Rightarrow t=2,\ t=\frac12$: \ $\res{x=1,\ x=-1}$

\ex{2} $x^2-1=8\Rightarrow x=\pm3$; \ $x=-3$ no vale: \ $\res{x=3}$

\ex{2} $x^2+1\gt0$ siempre, así que solo hace falta $x^2-4\le0$: \ $\res{x\in [-2,\, 2]}$\\[2pt]\begin{tikzpicture}[x=0.420cm,baseline=-3pt]
\draw[-latex,muted] (-3.6,0)--(3.7,0);
\draw[muted] (-3,0.1)--(-3,-0.1) node[below=-1pt,font=\scriptsize,text=muted]{$-3$};
\draw[muted] (-2,0.1)--(-2,-0.1) node[below=-1pt,font=\scriptsize,text=muted]{$-2$};
\draw[muted] (-1,0.1)--(-1,-0.1) node[below=-1pt,font=\scriptsize,text=muted]{$-1$};
\draw[muted] (0,0.1)--(0,-0.1) node[below=-1pt,font=\scriptsize,text=muted]{$0$};
\draw[muted] (1,0.1)--(1,-0.1) node[below=-1pt,font=\scriptsize,text=muted]{$1$};
\draw[muted] (2,0.1)--(2,-0.1) node[below=-1pt,font=\scriptsize,text=muted]{$2$};
\draw[muted] (3,0.1)--(3,-0.1) node[below=-1pt,font=\scriptsize,text=muted]{$3$};
\draw[accent,line width=2.6pt] (-2,0)--(2,0);
\fill[accent] (-2,0) circle (2.3pt);
\fill[accent] (2,0) circle (2.3pt);
\end{tikzpicture}

\ex{2} $5000\cdot1,04^t\ge8000\Rightarrow t\ge\dfrac{\log1,6}{\log1,04}\approx11,98$: \ $\res{12\ \text{años}}$


\end{document}
